\section{Conclusion}
\label{sec:conclusion}

We introduced \textbf{LUNA}, a self-supervised foundation model designed to address the challenge of topological heterogeneity in EEG analysis. By leveraging learned queries and cross-attention, LUNA unifies recordings with diverse electrode layouts into a fixed latent space, enabling montage-agnostic modeling. Through extensive experiments across abnormality detection, artifact recognition, slowing classification, and emotion recognition, we demonstrate that LUNA matches or surpasses state-of-the-art performance while offering substantial efficiency gains in FLOPs and memory usage. Critically, these benefits hold across all evaluated electrode configurations.

While LUNA achieves strong results, especially on heterogeneous montages, our analysis also reveals limitations. Performance on SEED-V suggests sensitivity to unseen channel topologies, likely stemming from reliance on positional encodings learned during pre-training. Addressing this limitation, through enhanced spatial generalization strategies or hybrid learned/geometric embeddings, is an important direction for future work.

More broadly, this work highlights the promise of topology-agnostic latent representations for scalable EEG modeling. Future extensions include exploring unified models across EEG and invasive modalities (e.g., sEEG, ECoG), integrating domain-specific priors (e.g., neurophysiological constraints), and adapting LUNA for real-time inference scenarios. Beyond technical advancements, the development of efficient, topology-invariant EEG models like LUNA could enhance neurological diagnostics and research accessibility. However, careful attention must be paid to mitigating risks such as algorithmic bias and ensuring patient data privacy for deployment. Future work should integrate ethical concerns alongside technical improvements. Pre-training montage diversity is limited (three dominant layouts); future work will incorporate multi-dataset pre-training and randomized channel dropout to improve generalization to unseen, dense and sparse montages.


